\documentclass[a4paper,14pt]{extarticle}

\usepackage[T2A]{fontenc}
\usepackage[utf8]{inputenc}
\usepackage[russian]{babel}
\usepackage[left=30mm,right=15mm,top=20mm,bottom=20mm]{geometry}
\usepackage{setspace}
\usepackage{indentfirst}
\usepackage{hyperref}

\onehalfspacing
\setlength{\parindent}{1.25cm}
\hypersetup{colorlinks=true, linkcolor=black, urlcolor=blue}

\begin{document}

\begin{titlepage}
  \begin{center}
    \textbf{Министерство науки и высшего образования Российской Федерации}\\[0.3cm]
    \textbf{Федеральное государственное автономное образовательное учреждение\\
    высшего образования}\\[0.3cm]
    \textbf{«НАЦИОНАЛЬНЫЙ ИССЛЕДОВАТЕЛЬСКИЙ УНИВЕРСИТЕТ ИТМО»}\\[0.5cm]
    \hrule
    \vspace{0.3cm}
    Факультет программной инженерии и компьютерной техники\\[0.2cm]
    Образовательная программа «Разработка программного обеспечения»\\[2cm]

    \large\textbf{ОТЧЁТ}\\[0.3cm]
    \textbf{по лабораторной работе №2 и домашнему заданию №5}\\[0.3cm]
    \textbf{по дисциплине «Бэкенд-разработка»}\\[0.5cm]
    \normalsize Тема: \textbf{Микросервисная архитектура, очереди сообщений}\\[3cm]

    \begin{flushright}
      \begin{tabular}{ll}
        Выполнил: & Иванов Виктор\\
        Группа:   & БР1.1\\
        Вариант:  & Система бронирования столиков\\
                  & в ресторанах\\[0.5cm]
        Проверил: & \\
      \end{tabular}
    \end{flushright}

    \vfill
    Санкт-Петербург, 2026
  \end{center}
\end{titlepage}

\tableofcontents
\newpage

\section{Цель работы}

Разделить монолитное приложение \texttt{lab1} на микросервисы с отдельными базами данных (принцип \textit{database-per-service}), описать межсервисные контракты в формате OpenAPI, реализовать асинхронное взаимодействие через RabbitMQ по образцу учебного репозитория.

\section{Выполненные задачи}

\subsection{ЛР2: микросервисы}

\begin{itemize}
  \item \textbf{user-service} (порт 4001): регистрация, аутентификация, JWT; внутренний метод \texttt{GET /internal/users/:id}.
  \item \textbf{restaurant-service} (порт 4002): каталог ресторанов и столиков; внутренний метод \texttt{GET /internal/tables/:tableId}; консюмер очереди RabbitMQ.
  \item \textbf{booking-service} (порт 4003): CRUD бронирований для владельца; перед созданием --- HTTP-вызовы к user и restaurant; проверка конфликтов по дате/времени в своей БД.
  \item \textbf{notification-service}: отдельный процесс, своя БД, запись логов по событиям из очереди.
\end{itemize}

Файлы SQLite: \texttt{data/user.sqlite}, \texttt{data/restaurant.sqlite}, \texttt{data/booking.sqlite}, \texttt{data/notification.sqlite}.

\subsection{ДЗ5: RabbitMQ}

Использован Docker-образ \texttt{rabbitmq:3.13-management-alpine}. После успешного создания брони \textbf{booking-service} публикует JSON в fanout-exchange \texttt{booking_events}. Подписчики: \textbf{restaurant-service} (таблица \texttt{booking_event_log}) и \textbf{notification-service} (таблица \texttt{notification_log}).

\section{Артефакты в репозитории}

\begin{itemize}
  \item \texttt{docs/design.md} --- технический дизайн и план миграции с монолита (соответствие ДЗ4).
  \item \texttt{openapi-inter-service.yaml} --- спецификация внутренних HTTP API.
  \item \texttt{README.md} --- порядок запуска и пример запроса.
  \item \texttt{docker-compose.yml} --- брокер сообщений.
\end{itemize}

\section{Выводы}

Монолитное API разбито на независимые сервисы с изолированными БД. Синхронная координация при создании брони выполняется через защищённые внутренние эндпоинты; уведомление смежных контекстов --- через очереди RabbitMQ, что снижает связность и позволяет масштабировать обработку событий отдельно от HTTP-потока.

\end{document}
